% Independent verification supplied for integration into the root E(c) statement.
\paragraph{The order-two tangential lift.}
Put $K=\mathbb Q(i)$ and $E:y^2=x^3-x$ with origin $O$. The local parameter $t=-x/y$ at $O$ satisfies $[-1]^*t=-t$. For the pro-$5$ punctured formal disc, take the compatible covers $t_n^{5^n}=t$, with $t_{n+1}^5=t_n$. The maps
\[
 t_n\longmapsto -t_n\qquad(n\geq0)
\]
commute with these transition maps, fix the coefficient field, and square to the identity. They are therefore compatible algebraic lifts of $t\mapsto-t$ on the entire pro-$5$ local tower. For a $K$-rational nonzero tangent vector $v$, they give the path from the tangential fibre at $v$ to the one at $-v$ whose root ratios are all $-1$. The path is Galois invariant, because these ratios belong to $K$, and its composite with its image under $[-1]$ is the identity path. Pullback of finite etale covers of $E\setminus\{O\}$ to this local tower therefore produces a genuine continuous automorphism $\alpha$ of the geometric pro-$5$ fundamental group based at $v$. It is an involution, and it commutes with the genuine action supplied by the tangential section $s$. This proves the lift on the specified quotient; a compatible path in the full profinite tower, including its $2$-power covers, is not used.

Let $N_2=\pi^{(5)}/\overline{[\pi^{(5)},[\pi^{(5)},\pi^{(5)}]]}$, with $V=N_2^{\rm ab}=T_5E$ and $Z=[N_2,N_2]\simeq\bigwedge^2V\simeq\mathbb Z_5(1)$. The group $N_2$ is the free pro-$5$ group of nilpotence class two on two generators. Its integral logarithmic coordinates require no division by $5$: its underlying $\mathbb Z_5$-module $L$ has addition
\[
 a\mathbin\oplus b=ab[a,b]^{-1/2},
\]
scalar multiplication $\lambda a=a^\lambda$, and bracket the group commutator. The class-two commutator identities prove the additive and bilinear identities directly; $1/2\in\mathbb Z_5$ makes every displayed operation defined and continuous. Equivalently its group law in $\mathbb Z_5^2\oplus\mathbb Z_5$ is
\[
 (u,z)(u',z')=(u+u',z+z'+\tfrac12(u\wedge u')).
\]
This model has the universal property of the free class-two pro-$5$ group: a word reduces to two generator powers and one central commutator power, and the displayed identities extend these powers continuously from integers to $\mathbb Z_5$.

On $V$, $\alpha$ is $-1$; on $Z=\bigwedge^2V$ it is $+1$. Consequently
\[
 L=L^-\oplus L^+,\qquad L^{\pm}=\tfrac12(1\pm\alpha)L,
 \qquad L^+=Z,\quad L^-\xrightarrow{\ \sim\ }V.
\]
The last equality follows because the projection of an $\alpha$-fixed vector to $V$ is killed by $2$, while $V$ is torsion free. Projection restricts to an isomorphism on $L^-$ because $Z$ is fixed and $2$ is a unit. Since the section action commutes with $\alpha$, it preserves this splitting. Thus its class-two action has no homomorphism $V\to Z$ added to the lift of its action on $V$.

The $K$-rational endomorphism $[i]:(x,y)\mapsto(-x,iy)$ commutes with Galois. Its polynomial $T^2+1$ splits into distinct linear factors over $\mathbb Z_5$; its two eigenspaces on $T_5E$ have rank one. Indeed its determinant is $1$ because it preserves the Weil pairing, and two identical eigenvalues would have determinant $-1$. The centralizer of $[i]$ on $V$ is therefore the abelian split Cartan $\mathbb Z_5^*\times\mathbb Z_5^*$. The representation of $\Gamma_K$ on $V$ factors through its abelianization, and $\Delta=\Gamma_{K^{\rm ab}}$ acts trivially on $V$. It also acts trivially on $Z=\bigwedge^2V$. The invariant splitting just proved yields
\[
 N_2^{s(\Delta)}=N_2.
\]
In particular the class-two centralizer contains every element with nonzero abelianization, while cusp inertia in this quotient is its central commutator subgroup. The equivariant quotient map $q:\pi^{(5)}\twoheadrightarrow N_2$ gives only
$q(C_{\pi^{(5)}}(s(\Delta)))\subset N_2^{s(\Delta)}$; no surjectivity of this map on fixed subgroups is implied or used. Lifting one of the noncentral invariant elements requires solving the successive nonabelian fixed-element equations through the lower central series.
